\chapter{应用实例}

\section{平面二连杆机械臂}

最简单的串联机器人是平面二连杆机构（2R 平面臂）。尽管结构简单，它包含 DH 参数法的所有基本元素。

\begin{figure}[htbp]
	\centering
	\begin{tikzpicture}[scale=2.0]
		% 基座
		\fill[pattern=north east lines] (-0.3,0) rectangle (0.3,0.2);
		
		% 关节0
		\draw[thick, fill=blue!20] (0,0) circle (0.08);
		\node[below left] at (0,0) {$O_0$};
		
		% 连杆1 (L1)
		\draw[very thick, red!60!black] (0,0) -- (1.2,0.6);
		\node[red!60!black] at (0.5,0.1) {$L_1$};
		
		% 关节1
		\draw[thick, fill=blue!20] (1.2,0.6) circle (0.08);
		\node[above right] at (1.2,0.6) {$O_1$};
		
		% theta1
		\draw[->] (0.4,0) arc (0:26:0.4);
		\node at (0.55,0.15) {$\theta_1$};
		
		% 连杆2 (L2)
		\draw[very thick, red!60!black] (1.2,0.6) -- (1.9,1.5);
		\node[red!60!black] at (1.7,0.9) {$L_2$};
		
		% theta2
		\draw[dashed] (1.2,0.6) -- (1.8,0.9);
		\draw[->] (1.6,0.8) arc (26:70:0.4);
		\node at (1.55,1.0) {$\theta_2$};
		
		% 末端
		\filldraw (1.9,1.5) circle (1.5pt) node[right] {$P$};
		
		% 坐标系标注
		\draw[->, blue] (0,0) -- (0.4,0) node[below] {$x_0$};
		\draw[->, blue] (0,0) -- (0,0.4) node[left] {$y_0$};
	\end{tikzpicture}
	\caption{平面二连杆机械臂（2R 平面臂）}
\end{figure}

\subsection{坐标系建立与 DH 参数表}

约定 $z_0$ 轴垂直纸面向外（$z$ 轴沿关节旋转轴）。$x_0$ 水平向右，$y_0$ 竖直向上。所有 $z_i$ 轴均垂直于纸面向外（平行）。

\begin{center}
	\begin{tabular}{ccccc}
		\toprule
		$i$ & $a_{i-1}$ & $\alpha_{i-1}$ & $d_i$ & $\theta_i$ \\
		\midrule
		1 & 0   & 0 & 0 & $\theta_1$ \\
		2 & $L_1$ & 0 & 0 & $\theta_2$ \\
		\bottomrule
	\end{tabular}
\end{center}

\subsection{变换矩阵计算}

\[
{}^{0}_{1}\mathbf{T} = \begin{bmatrix}
	\cos\theta_1 & -\sin\theta_1 & 0 & 0 \\
	\sin\theta_1 &  \cos\theta_1 & 0 & 0 \\
	0            &  0            & 1 & 0 \\
	0            &  0            & 0 & 1
\end{bmatrix}
\]

\[
{}^{1}_{2}\mathbf{T} = \begin{bmatrix}
	\cos\theta_2 & -\sin\theta_2 & 0 & L_1 \\
	\sin\theta_2 &  \cos\theta_2 & 0 & 0 \\
	0            &  0            & 1 & 0 \\
	0            &  0            & 0 & 1
\end{bmatrix}
\]

\[
{}^{0}_{2}\mathbf{T} = {}^{0}_{1}\mathbf{T} \cdot {}^{1}_{2}\mathbf{T} = \begin{bmatrix}
	\cos(\theta_1+\theta_2) & -\sin(\theta_1+\theta_2) & 0 & L_1\cos\theta_1 \\
	\sin(\theta_1+\theta_2) &  \cos(\theta_1+\theta_2) & 0 & L_1\sin\theta_1 \\
	0                        &  0                        & 1 & 0 \\
	0                        &  0                        & 0 & 1
\end{bmatrix}
\]

\textbf{末端位置：}$x = L_1\cos\theta_1$，$y = L_1\sin\theta_1$ —— 等等，这不对。让我们重新检查。

实际上平面 2R 臂的连杆参数应当这样取：

如果采用\textbf{改进 DH}（Craig 约定），并约定 $x_0$ 沿连杆 1 方向（$\theta_1=0$ 时），$y_0$ 与之垂直：

设 $z_i$ 全部指向纸外。连杆 1 长度 $l_1$，连杆 2 长度 $l_2$。

\begin{center}
	\begin{tabular}{ccccc}
		\toprule
		$i$ & $a_i$ & $\alpha_i$ & $d_i$ & $\theta_i$ \\
		\midrule
		1 & 0    & 0 & 0      & $\theta_1$ \\
		2 & $l_1$ & 0 & 0      & $\theta_2$ \\
		3 & $l_2$ & 0 & 0      & 0 \\
		\bottomrule
	\end{tabular}
\end{center}

其中 $\{3\}$ 为工具坐标系（位于末端）。

计算：
\[
{}^{0}_{3}\mathbf{T} = \begin{bmatrix}
	\cos(\theta_1+\theta_2) & -\sin(\theta_1+\theta_2) & 0 & l_1\cos\theta_1 + l_2\cos(\theta_1+\theta_2) \\
	\sin(\theta_1+\theta_2) &  \cos(\theta_1+\theta_2) & 0 & l_1\sin\theta_1 + l_2\sin(\theta_1+\theta_2) \\
	0                        &  0                        & 1 & 0 \\
	0                        &  0                        & 0 & 1
\end{bmatrix}
\]

末端位置 $\mathbf{p} = (l_1\cos\theta_1 + l_2\cos(\theta_1+\theta_2),\, l_1\sin\theta_1 + l_2\sin(\theta_1+\theta_2),\, 0)\T$，与直觉一致。

\section{SCARA 机器人}

SCARA（Selective Compliance Assembly Robot Arm）是一种广泛应用于装配作业的四自由度机器人。它有三个旋转关节和一个移动关节。

\begin{figure}[htbp]
	\centering
	\begin{tikzpicture}[scale=1.3]
		% 基座
		\fill[pattern=north east lines, pattern color=gray] (-0.5,0) rectangle (0.5,0.3);
		\node at (0,0.15) {基座};
		
		% 关节1
		\draw[very thick] (0,0.3) -- (0,0.6);
		\node[left] at (0,0.45) {$J_1$};
		
		% 连杆1 (水平)
		\draw[very thick, red!60!black] (0,0.6) -- (2.5,0.6);
		\node[above] at (1.25,0.62) {$L_1$};
		
		% 关节2
		\node[right] at (2.5,0.6) {$J_2$};
		
		% 连杆2 (水平)
		\draw[very thick, red!60!black] (2.5,0.6) -- (4.2,0.9);
		\node[above] at (3.35,0.77) {$L_2$};
		
		% 关节3 (移动关节, 垂直)
		\draw[very thick, blue] (4.2,0.9) -- (4.2,2.0);
		\node[right] at (4.2,1.45) {$J_3$ (P)};
		\draw[<->] (4.5,0.9) -- (4.5,2.0) node[midway, right] {$d_3$};
		
		% 关节4 (旋转, 末端)
		\draw[very thick, blue] (4.2,2.0) -- (4.2,2.3);
		\node[right] at (4.2,2.15) {$J_4$};
		
		% 末端
		\draw[fill=orange!30] (3.9,2.3) rectangle (4.5,2.5);
		\node at (4.2,2.4) {末端};
		
		% theta1, theta2
		\draw[->] (0.5,0.6) arc (0:0:0.5);
		\node at (1.0,0.3) {$\theta_1$};
		\draw[->] (2.8,0.58) arc (0:7:0.4);
		\node at (3.1,0.5) {$\theta_2$};
	\end{tikzpicture}
	\caption{SCARA 机器人结构示意图}
\end{figure}

\subsection{DH 参数表（改进 DH 约定）}

约定 $z_0$ 垂直向上。各关节轴方向：
\begin{itemize}
	\item $z_0$：垂直向上（沿 $J_1$ 轴）
	\item $z_1$：垂直向上（沿 $J_2$ 轴）
	\item $z_2$：垂直向下（沿 $J_3$ 的移动方向）
	\item $z_3$：垂直向下（沿 $J_4$ 轴）
\end{itemize}

\begin{center}
	\begin{tabular}{cccccc}
		\toprule
		$i$ & $a_i$ & $\alpha_i$ & $d_i$ & $\theta_i$ & 关节 \\
		\midrule
		1 & $l_1$ & 0          & $d_1$   & $\theta_1$ & R \\
		2 & $l_2$ & $180^\circ$ & 0       & $\theta_2$ & R \\
		3 & 0     & 0          & $-d_3^*$ & 0          & P \\
		4 & 0     & 0          & $d_4$   & $\theta_4$ & R \\
		\bottomrule
	\end{tabular}
\end{center}

\textit{* $d_3$ 取负是因为 $z_2$ 向下，正值对应末端下降。}

\section{标准 DH 示例：PUMA 560}

PUMA 560 是机器人学中最经典的六自由度工业机器人。此处给出其标准 DH 参数，供实践参考。

\begin{center}
	\begin{tabular}{cccccc}
		\toprule
		$i$ & $a_{i-1}$ & $\alpha_{i-1}$ & $d_i$ & $\theta_i$ \\
		\midrule
		1 & 0          & 0          & 0           & $\theta_1$ \\
		2 & 0          & $-90^\circ$ & $d_2$       & $\theta_2$ \\
		3 & $a_2$      & 0          & 0           & $\theta_3$ \\
		4 & $a_3$      & $-90^\circ$ & $d_4$       & $\theta_4$ \\
		5 & 0          & $90^\circ$  & 0           & $\theta_5$ \\
		6 & 0          & $-90^\circ$ & 0           & $\theta_6$ \\
		\bottomrule
	\end{tabular}
\end{center}

典型数值：$d_2 = 149.09\text{ mm}$，$a_2 = 431.8\text{ mm}$，$a_3 = -20.32\text{ mm}$，$d_4 = 433.07\text{ mm}$。

\subsection{写出变换矩阵}

对旋转关节，$\theta_i$ 是变量。将 DH 参数依次代入标准 DH 矩阵公式即可得到每个 ${}^{i-1}_{i}\mathbf{T}$。整体正运动学为：

\[
{}^{0}_{6}\mathbf{T} = {}^{0}_{1}\mathbf{T}(\theta_1) \cdot {}^{1}_{2}\mathbf{T}(\theta_2) \cdot {}^{2}_{3}\mathbf{T}(\theta_3) \cdot {}^{3}_{4}\mathbf{T}(\theta_4) \cdot {}^{4}_{5}\mathbf{T}(\theta_5) \cdot {}^{5}_{6}\mathbf{T}(\theta_6)
\]

前三个关节决定腕心位置（称为\textbf{定位机构}），后三个关节决定末端姿态（称为\textbf{定向机构}）。这种"位置-姿态解耦"是六自由度腕部偏置型机器人的重要特征。

\section{DH 参数法的局限与注意事项}

\begin{myTheorem}{DH 参数法的局限性}
	\begin{enumerate}
		\item \textbf{非唯一性}：DH 参数表不唯一。坐标系原点沿 $z$ 轴平移会改变 $d$，$x$ 轴的可选方向也能导致多组参数。
		\item \textbf{平行关节轴的歧义}：当相邻关节轴平行时，$a=0$ 可选与否（取公垂线与前一公垂线相交或不相交均可）导致 $d$ 值不同。
		\item \textbf{不能描述闭链}：DH 参数法仅适用于串联开链机构。并联机器人或含闭链的机构需要其他方法。
		\item \textbf{数值敏感性}：当相邻关节轴接近平行时，微小的制造误差可导致 DH 参数剧烈变化（奇异性问题）。
	\end{enumerate}
\end{myTheorem}

针对这些局限，机器人学界发展出了多种替代方案，如\textbf{螺旋轴参数法}（Product of Exponentials, PoE）、\textbf{Hayati 参数}（针对平行轴情况）等。但对于标准串联机器人而言，DH 参数法仍然是教学和工程实践中最主流的方法。

\section{小结}

DH 参数法将多自由度串联机器人的运动学建模简化为以下步骤：

\[
\boxed{\text{画出关节轴}} \;\to\; \boxed{\text{建立连杆坐标系}} \;\to\; \boxed{\text{提取 DH 参数}} \;\to\; \boxed{\text{写出变换矩阵}} \;\to\; \boxed{\text{连乘得末端位姿}}
\]

这一方法以其\textbf{系统性、简洁性}和\textbf{可编程性}，成为现代机器人学的基石。熟练掌握 DH 参数法，是进入机器人运动学、动力学和控制领域的第一步。
